\clearpage\thispagestyle{empty}
\pdfbookmark[0]{Quellen zur Besprechung von Nouets astronomischen Beobachtungen in Ägypten}{nouet1811-sources}
\begin{center}\textbf{Quellen zur Besprechung von Nouets astronomischen Beobachtungen in Ägypten}\end{center}
\noindent Diese KI-übersetzte Fassung in modernem Deutsch folgt Carl Friedrich Gauß, \textit{Werke}, Band~VI (Göttingen, 1874), S.~534. Die Werkausgabe enthält einen Auszug aus einer längeren Besprechung in den \textit{Göttingischen gelehrten Anzeigen}, Stück~79, 18. Mai 1811, S.~780--781.
\par\smallskip\noindent Die Besprechung gilt Nouets \textit{Observations astronomiques faites en Egypte pendant les années VI, VII et VIII}, in der \textit{Description de l'Egypte}, \textit{État moderne}, Band~I, erste Lieferung (Paris, Imprimerie impériale, 1809), S.~1--20. Die \href{https://archive.org/details/bibliographicala00thom}{Bibliografie der London Institution}, S.~26--27, verzeichnet die Abhandlung. Der \href{https://catalogue.bnf.fr/ark:/12148/cb35129066r}{Katalog der Bibliothèque nationale de France} weist den Band nach.
\par\smallskip\noindent Im Erstdruck ist die letzte Ziffer von 3430 klar als Null zu erkennen. Die Fassung übernimmt aus der Werkausgabe die Schreibung Ssalehhiyeh und den Punkt vor der Fortsetzung des Satzes über die magnetische Deklination.
\par\smallskip\noindent \href{https://gdz.sub.uni-goettingen.de/id/PPN236007467}{Text der Werkausgabe}. Erstdruck, \href{https://images.p.sub.uni-goettingen.de/iiif/image/gdz:PPN319721507_1811:00000812/full/1650,/0/default.jpg}{S.~780} und \href{https://images.p.sub.uni-goettingen.de/iiif/image/gdz:PPN319721507_1811:00000813/full/1650,/0/default.jpg}{S.~781}.
\par\smallskip\noindent Übersetzung und Satz mit KI erstellt.
\clearpage
